
d'une
démonstration de la formule analogue dans le cas analytique complexe (ou rigide-analytique !). Rappelons que dans le cas analytique complexe, seul le cas d'un faisceau localement libre @@M0@@ sur une variété non singulière compacte @@M1@@, au-dessus d'un @@M2@@ ponctuel, est actuellement connu (grâce à la formule de l'index de Atiyah--Singer), le cas où @@M3@@ est seulement supposé cohérent semblant échapper encore aux techniques connues jusqu'à présent.

\textbf{2.2.} Comme ingrédient essentiel de la formule de Riemann--Roch, rappelons aussi pour mémoire le \uline{problème de finitude} soulevé dans Exp. III: si @@M4@@ est un morphisme propre, et @@M5@@ un complexe sur @@M6@@ qui est pseudo-cohérent relativement à @@M7@@, alors @@M8@@ est-il pseudo-cohérent? Une démonstration générale de cette conjecture mettra probablement en oeuvre des idées algébriques qui permettront de donner une démonstration du ``théorème'' de finitude de Grauert [7] sous des conditions notablement plus générales que dans loc. cit., en se plaçant d'emblée sur des espaces (ou topos) de base topologiquement annelés en @@M9@@-algèbres @@M10@@ plus généraux que des espaces analytiques (englobant le cas où @@M11@@ serait un espace localement compact avec le faisceau des fonctions continues complexes, ou une variété @@M12@@ avec le faisceau des fonctions complexes @@M13@@, ou un espace analytique réel, etc).

\textbf{2.3.} Mentionnons enfin qu'il y aurait lieu de remplacer dans ces questions l'hypothèse de propreté sur le morphisme @@M14@@ par l'hypothèse que les faisceaux de cohomologie @@M15@@ des complexes envisagés sur @@M16@@ aient leurs supports propres sur @@M17@@.

\section*{3. Formule de Riemann--Roch ``sans dénominateurs'' pour une immersion}

\textbf{3.1.} Soit @@M18@@ une immersion fermée régulière (Exp. VII), de faisceau conormal @@M19@@. Supposons pour simplifier que @@M20@@ admet un faisceau inversible ample. La formule de Riemann--Roch pour @@M21@@~:
@@M22@@
jointe à la formule générale (Exp. VII)~:
@@M23@@
(en supposant @@M24@@ de codimension @@M25@@), permet d'exprimer les classes de Chern de @@M26@@ dans @@M27@@ sous la forme
